\mxnewpage\section*{Correction EV2-A}\label{EV2-A-CORR}
\textbf{QCM :} 1 C, 2 D, 3 B, 4 C, 5 B, 6 C, 7 D, 8 B, 9 C, 10 D, 11 B, 12 D. Chaque réponse correcte : \num{0,5} point ; total 6.
\paragraph{Vérifications des automatismes.} Priorités ; somme de fractions ; soustraction des exposants ; coefficient multiplicateur \num{1,2} ; division par \num{0,8} ; produit \num{1,25} fois \num{0,8} ; résolution affine ; accroissement vertical divisé par l’horizontal ; moyenne ; angle divisé par 360 ; complément à 1 ; substitution entre parenthèses.
\subsection*{Exercice 1 -- Une probabilité inconnue}
\begin{enumerate}
\item $\dfrac{3}{10}$. \textit{Barème : 1 point. Résultat et justification élémentaire attendus ; demi-point possible pour démarche correcte inachevée.}
\item $P(\overline A\cap B)=\dfrac{1}{2}-\dfrac{3}{10}=\dfrac{1}{5}$ ; diviser par $\dfrac{3}{5}$ donne $\dfrac{1}{3}$ ; compléments $\dfrac{1}{4}$ et $\dfrac{2}{3}$. \textit{Barème : 2 points. Méthode ou justification : 1 ; calcul et conclusion : 1.}
\item $P_B(A)=\dfrac{3}{5}$ ; différent de $P(A)=\dfrac{2}{5}$, donc non indépendants. \textit{Barème : 2 points. Méthode ou justification : 1 ; calcul et conclusion : 1.}
\end{enumerate}
\subsection*{Exercice 2 -- Position relative et variations}
\begin{enumerate}
\item $f\prime=3x(x-2)$ ; croît jusqu’à 0, décroît jusqu’à 2, croît ensuite. \textit{Barème : 1 point. Résultat et justification élémentaire attendus ; demi-point possible pour démarche correcte inachevée.}
\item $f-g=x^2(x-3)$ : négatif pour x<3 sauf égalité en 0, nul aussi en 3, positif après 3. \textit{Barème : 2 points. Méthode ou justification : 1 ; calcul et conclusion : 1.}
\item $f(-1)=-2,f(0)=2,f(2)=-2,f(4)=18$. Minimum -2 en -1 et 2 ; maximum 18 en 4. \textit{Barème : 2 points. Méthode ou justification : 1 ; calcul et conclusion : 1.}
\end{enumerate}
\subsection*{Exercice 3 -- Une médiatrice}
\begin{enumerate}
\item I(3;3), normale (6;2), ou (3;1). \textit{Barème : 1 point. Résultat et justification élémentaire attendus ; demi-point possible pour démarche correcte inachevée.}
\item $3x+y-12=0$. \textit{Barème : 1 point. Résultat et justification élémentaire attendus ; demi-point possible pour démarche correcte inachevée.}
\item $3(2)+6-12=0$ donc oui ; $PA^2=4+16=20$, $PB^2=16+4=20$. \textit{Barème : 2 points. Méthode ou justification : 1 ; calcul et conclusion : 1.}
\end{enumerate}
\paragraph{Étalonnage.} Durée cible : 20–25 min pour les automatismes, 80–85 min pour les exercices, 10–15 min de relecture. Ces durées sont des estimations éditoriales à mesurer en classe. Accepter toute méthode correcte ; ne pas sanctionner deux fois une erreur numérique propagée avec cohérence. Le barème proposé ne vaut pas décision d’harmonisation de l’équipe.
